\slajdid{09_3-s033}
\begin{frame}[label=09_3-s033]
\begin{columns}[c,onlytextwidth]\column{.49\textwidth}\centering\input{slike/tikz/s033_pdn_mreza.tex}\column{.49\textwidth}\centering\begin{tikzpicture}[x=.7cm,y=.58cm,font=\sitneoznake,>=Latex]\begin{scope}[shift={(0,3.25)}]\draw(0,-.35)--(0,.35);\draw(-.13,-.4)--(-.13,.4);\draw(0,.3)--(.4,.3)--(.4,.65);\draw(0,-.3)--(.4,-.3)--(.4,-.65);\draw(-.23,0)circle(.1);\draw(-1,0)node[left]{$\overline A$}--(-.33,0);\end{scope}\begin{scope}[shift={(3.2,4)}]\draw(0,-.35)--(0,.35);\draw(-.13,-.4)--(-.13,.4);\draw(0,.3)--(.4,.3)--(.4,.65);\draw(0,-.3)--(.4,-.3)--(.4,-.65);\draw(-.23,0)circle(.1);\draw(-1,0)node[left]{B}--(-.33,0);\end{scope}\begin{scope}[shift={(3.2,2.5)}]\draw(0,-.35)--(0,.35);\draw(-.13,-.4)--(-.13,.4);\draw(0,.3)--(.4,.3)--(.4,.65);\draw(0,-.3)--(.4,-.3)--(.4,-.65);\draw(-.23,0)circle(.1);\draw(-1,0)node[left]{C}--(-.33,0);\end{scope}\begin{scope}[shift={(1.6,0.6)}]\draw(0,-.35)--(0,.35);\draw(-.13,-.4)--(-.13,.4);\draw(0,.3)--(.4,.3)--(.4,.65);\draw(0,-.3)--(.4,-.3)--(.4,-.65);\draw(-.23,0)circle(.1);\draw(-1,0)node[left]{D}--(-.33,0);\end{scope}\draw(.4,3.9)--(.4,5)--(3.6,5)--(3.6,4.65);\draw(3.6,5)--(3.6,5.3);\draw(3.2,5.3)--(4,5.3)node[above left]{$V_{DD}$};\draw(3.6,3.35)--(3.6,3.15);\draw(.4,2.6)--(.4,1.5)--(3.6,1.5)--(3.6,1.85);\draw(2,1.5)--(2,1.25);\draw(2,-.05)--(2,-.4);\node[below]at(2,-.7){PUN mreža};\end{tikzpicture}
\end{columns}\par\bigskip
Prilikom pravljenja PUN mreže mogli smo odmah da uzmemo u obzir prirodu dualnosti ovih mreža. Posmatrajući PDN mrežu - paralelne veze iz PDN mreže postaju redne veze u PUN mreži, redne veze iz PDN mreže postaju paralelne veze u PUN mreži. Na primer: tranzistor sa ulazom D je bio u PDN mreži paralelno povezan sa ostatkom mreže, dok će u PUN mreži biti redno povezan sa ostatkom mreže. Itd…
\end{frame}
